% =====================================================================
%  Attributing RoPE: exact score attribution, and what actually limits it
%
%  Self-contained NeurIPS-format preprint (two-column review layout).
%  Build:  pdflatex main && bibtex main && pdflatex main && pdflatex main
%
%  EVERY quantitative claim in this paper is traceable to one of:
%    results/measurements.json   (canonical, committed)
%    figures/figNN_*.csv          (CSV sidecar of the plotted array)
%  The appendix section "Provenance of every number" carries the audit.
%
%  Figures are included by relative path and guarded with \IfFileExists
%  so that the document still compiles with an empty ../figures/.
% =====================================================================
\documentclass[twocolumn]{article}

% ---- review/preprint setup, self-contained (no .sty files required) ----
\usepackage{natbib}
% Geometry: give the package ONE non-contradictory specification. The
% previous pair (\usepackage[margin=1in]{geometry} followed by
% \geometry{textwidth=5.5in,textheight=9in}) made geometry emit
%   "Over-specification in `h'-direction. `width' (397.485pt) is ignored"
% and lay the document out at 6.5in wide, silently discarding the declared
% 5.5in. textheight=9in on letterpaper reproduces the intended 1in top and
% bottom margins; the 5.5in block is then centred (1.5in left/right).
\usepackage[letterpaper,textwidth=5.5in,textheight=9in]{geometry}
\setlength{\columnsep}{0.25in}
\usepackage{amsmath,amssymb,amsthm,bm}
\usepackage{graphicx}
\usepackage{booktabs}

% natbib already defines \bibsep as a *length* (it is used as
% \setlength{\itemsep}{\bibsep} inside thebibliography), so widen that
% register rather than trying to \newcommand it.
\renewcommand{\bibname}{References}
\setlength{\bibsep}{\bigskipamount}

% ---- theorem environments ----
\theoremstyle{plain}
\newtheorem{theorem}{Theorem}
\newtheorem{proposition}[theorem]{Proposition}
\theoremstyle{remark}
\newtheorem{remark}[theorem]{Remark}

% ---- figure guard: include if present, otherwise typeset a labelled box ----
% \paperfigure{<file stem under ../figures/>}{<caption>}{<label>}
\newcommand{\paperfigure}[3]{%
  \begin{figure}[t]
    \centering
    \IfFileExists{../figures/#1.png}{%
      % --- true branch: the PNG is on disk -------------------------
      \includegraphics[width=\linewidth]{../figures/#1.png}%
    }{%
      % --- false branch: typeset a labelled placeholder -------------
      \fbox{%
        \parbox[c][3.2cm][c]{0.94\linewidth}{%
          \centering
          {\footnotesize\ttfamily ../figures/#1.png}\\[4pt]
          {\footnotesize (figure file absent at build time; the CSV
           sidecar in ../figures/ carries the plotted array)}%
        }%
      }%
    }%
    \caption{#2}%
    \label{#3}%
  \end{figure}%
}%

% ---- inline helpers ----
\newcommand{\rope}{\textsc{rope}}
\newcommand{\yarn}{\textsc{yarn}}
\newcommand{\pip}{\textsc{pp-rope}}
% short handles used in Related Work only
\newcommand{\alpibib}{\textsc{alibi}}
\newcommand{\lengthxtrap}{\textsc{length-extrap.}}
\newcommand{\longrope}{\textsc{longrope}}
\newcommand{\ropetheory}{\textsc{round-and-round}}
\newcommand{\R}{R}
\newcommand{\invf}{\mathrm{inv\_freq}}
\newcommand{\D}{\mathrm{D}_k}
\newcommand{\Ak}{A_k}
\newcommand{\Bk}{B_k}
\newcommand{\Rk}{R_k}
\newcommand{\psik}{\psi_k}
\newcommand{\half}{d/2}
\newcommand{\rad}{\mathrm{rad}}

\title{\vspace{-2.5ex}%
  \textbf{Attributing Rotary Position Embeddings: Exact Score Attribution,
  and What Actually Limits It}%
  \vspace{-1.5ex}}

\author{Ilya Slyatski}

\begin{document}

% With the twocolumn class option, article's \maketitle already spans both
% columns; do not wrap it in another \twocolumn.
\maketitle

\begin{abstract}
Rotary position embeddings (RoPE) \citep{su2021ropeformer} are the default
positional encoding in most modern transformers, and interpretability work now
routinely attributes attention scores to learned features using an additive
decomposition. We take the structure of that decomposition seriously and ask a
narrow question: \emph{exactly} what does RoPE do to it? Three things. First,
for a fixed pair of positions, the RoPE attention score is exactly bilinear in
the content vectors, so the additive feature decomposition
$\mathrm{score} = \sum_{ij} f_i g_j A_{ij}(\delta)$ is exact, not approximate;
we measure the residual at $3.6\times10^{-14}$. Because the rotary transform is
orthogonal, the magnitude gate $|q|$ is exactly position-independent, which
promotes the gate/phase split from heuristic to theorem. Second, the exactness
does not extend to position-\emph{free} coefficients: $A_{ij}$ is a function of
the relative distance $\delta$, and we measure one feature's contribution
reversing sign across the distance range while
the total remains additive to $5.0\times10^{-14}$. A position-free per-feature
attribution number is therefore not well-posed under RoPE; the per-feature,
per-distance one is. Third, we measure what context extension actually buys.
YaRN \citep{peng2023yarn} keeps $\mathrm{base} = 10000$ and ramps per frequency;
consequently $\max_k|\D|$ is bit-identical to plain RoPE at every one of $54$
tested distances, while the median angle falls by a constant factor of $0.4464$
and the linearizable fraction rises. Position interpolation
\citep{chen2023positionalinterp} and the common ``raise the base'' heuristic are
measured as controls. The honest conclusion is that context extension buys a
larger linearizable \emph{majority}, not a safe global linearization of the
head, because a few fast channels dominate the error. Partial RoPE is the clean
exception: at $p = 0.25$ the unrotated sub-score is position-free to exactly
zero spread. All numbers here are computed from a from-scratch implementation;
this is a structural characterization, and we state explicitly what it does not
include.
\end{abstract}

% =====================================================================
\section{Introduction}

Feature-level attribution has become a working tool in mechanistic
interpretability. Given a decomposition of an activation into a sum of learned
features, one asks how much each feature contributed to a circuit's output.
Applied to attention, this runs into a structural question that is usually
glossed over: attention scores are computed between a query at one position and
a key at another, and under a positional encoding, ``how much did feature $i$
contribute'' is not obviously a scalar.

Under RoPE it is not a scalar, and the reason is specific. In this paper we
derive the score in the relative frame, where RoPE's orthogonality makes the
query unrotated and the key rotated by $\R_{\delta}$, $\delta = n - m$.
Expanding that inner product channel by channel yields an exact closed form: the
head decomposes into $\half$ rotary pairs, each contributing a single sinusoid
in $\delta$ whose amplitude is fixed by the content vectors and whose phase
moves linearly with distance. The amplitude is a \emph{gate}; the phase is the
whole of the position dependence. This is a cleaner object than the score
itself, and it is exactly computable.

Three results follow, and we state them as precisely as we can.

\paragraph{R1: attribution is exact, and the gate/phase split is a theorem.}
For a fixed position pair the score is exactly bilinear in $(q, k)$, so the
additive decomposition holds to $3.6\times10^{-14}$. Moreover, because
$\R_{\delta}$ is orthogonal, $|q|$ and $|k|$ are exactly
position-independent. Any claim that RoPE corrupts bilinearity, or smears the
gate with position, is false at the score level, and we measure it to be false
rather than asserting it.

\paragraph{R2: the real obstacle is position-conditionality, not bilinearity.}
$A_{ij}$ depends on $\delta$. A single feature's contribution is a function of
distance, not a number. We measure that every feature's contribution reverses sign across a
$54$-point distance grid, against an additivity residual of
$5.0\times10^{-14}$. The well-posed object under RoPE is a per-feature,
per-distance attribution. A position-free scalar per feature is not.

\paragraph{R3: context extension buys a majority, not a global linearization.}
The common story is that RoPE's angles grow without bound and methods fix it by
damping them. Measuring $\max_k |\D|$, $\operatorname{median}_k |\D|$ and the
fraction of channels admitting $\cos \D \mapsto 1 - \D^2/2$,
$\sin \D \mapsto \D$ across RoPE, \yarn, position interpolation and the
``$\mathrm{base} = 500000$'' heuristic, we find $\max_k |\D|$ is
\emph{identical} for RoPE and \yarn\ at every distance tested. \yarn\
deliberately leaves the fastest channel alone. It lowers the median and raises
the linearizable fraction, but the amplitude-weighted error ratio
\yarn/RoPE is $0.99684$ at $\delta = 8192$. The fast channels dominate the
error and \yarn\ does not touch them.

\subsection{What we do not claim}

This matters enough to state before the results. The algebraic claims are made on
synthetic query and key vectors from a fixed seed, projected through synthetic
feature directions; Section~\ref{sec:trained} then repeats the central
measurements on activations from three trained checkpoints, and that section
reports what those activations do and do not settle. There is no sparse
autoencoder anywhere in this work: we attribute over a synthetic feature set,
with $8$ random directions and non-negative coefficients, which mimic the
\emph{shape} of a sparse code and are a test fixture for the algebra, not a model
of a real feature set.
There is no retrieval benchmark and no accuracy measurement: we report no
perplexity, no needle-in-a-haystack score and no long-context retrieval result.
We cannot and do not report anything of the form ``the model retrieves the
needle with accuracy $0.7$'', because nothing like that is measured. Where we
compare RoPE and \yarn\ we compare \emph{quantities of the position encoding}
--- angles, linearizable fractions, error ratios. We do not claim any scheme is
better in practice: all are in real use, our table ranks them on angle geometry
only, and Section~\ref{sec:context} shows that a favourable ranking on one axis
does not transfer to the amplitude-weighted one. The conclusions are
model-independent by construction, which is deliberate: they are claims about
the structure of the score under a particular positional encoding, and they
would hold for any head using it.

The contribution, stated exactly: a correct, carefully measured characterization
of what RoPE does to feature-level attention attribution, an identification of
position-conditionality as the operative obstruction, and a measurement of the
actual effect of context extension on the quantities that determine whether a
head can be linearized.

% =====================================================================
\section{Background}
\label{sec:background}

\subsection{Rotary position embeddings}

Let $d$ be an even head dimension and let $x_m \in \mathbb{R}^d$ be the
pre-projection content at absolute position $m$. RoPE defines
\begin{equation}
  \begin{split}
    \invf_k &= \mathrm{base}^{-2k/d},\; k = 0, \dots, \half - 1,\\
    \mathrm{base} &= 10000 .
  \end{split}
\end{equation}
and applies a block-diagonal rotation $\R_m$ whose $k$-th $2\times2$ block
rotates by the angle $m \cdot \invf_k$. With $q = W_Q x_m$ and $k = W_K x_n$ the
projections taken before rotation, the score at query position $m$ and key
position $n$ is
\begin{equation}
  s(m,n) \;=\; \bigl(\R_m q\bigr)^{\!\top} \bigl(\R_n k\bigr).
  \label{eq:rope}
\end{equation}
This is the formulation of \citet{su2021ropeformer}. Our implementation follows
the split-half convention, pairing channel $i$ with $i + \half$, so that one
frequency vector of length $\half$ drives all $d$ channels.

\subsection{The relative frame}

The rotation matrices $\R_m$ are orthogonal, hence $\R_m^{\!\top}\R_m = I$.
Factoring through the relative frame gives
\begin{equation}
  \begin{split}
    \R_m^{\!\top} \bigl(\R_n k\bigr) \;=\; \R_m^{\!\top}\R_n k \;=\; \R_{n-m} k,\\
    \R_m^{\!\top} \bigl(\R_m q\bigr) \;=\; q .
  \end{split}
\end{equation}
and therefore
\begin{equation}
  s(m,n) \;=\; q^{\!\top} \R_{\delta} k, \qquad \delta = n - m,
  \label{eq:relative}
\end{equation}
which depends on the two absolute positions only through their difference,
reducing the score of Equation~\ref{eq:rope} to a function of one number. This is
the relative-position property, and it is what makes the rest of the paper
possible: from here on, ``position'' means $\delta$ and nothing else.

\begin{proposition}[The gate is position-free]
\label{prop:gate}
$\bigl|\R_{\delta} k\bigr| = \bigl|k\bigr|$ for all $\delta$. Consequently the
magnitude of the query and of the key is exactly position-independent, and all
position information in $s(\delta)$ is carried by the rotation angle alone.
\end{proposition}

Proposition~\ref{prop:gate} is the load-bearing statement for the gate/phase
split used in sparse-feature interpretability: the ``gate'' $|q|$ is not an
approximation of the position-free part, it is \emph{exactly} the
position-free part, and the ``phase'' is the entire remainder. We measure the
numerical residual of this orthogonality in Section~\ref{sec:exact}.

\subsection{The per-pair closed form}
\label{sec:pair}

Group the head into rotary pairs $(k, k + \half)$, so that $\R_{\delta}$ mixes
exactly these two channels. In the split-half convention, the $k$-th block of
$\R_{\delta}$ acts on $(q[k], q[k+\half])$ and $(k[k], k[k+\half])$ as
\begin{equation}
  \begin{split}
    \R_\delta^{(k)} \;=\;
    \begin{pmatrix} \cos \D & -\sin \D \\ \sin \D & \cos \D \end{pmatrix},
    \\[-1pt]
    \qquad \D = \delta \invf_k .
  \end{split}
\end{equation}
with the same $\D$ on both diagonals. Expanding $q^{\!\top}\R_{\delta}k$ over
that block and collecting terms gives, exactly,
\begin{align}
  \Ak &= q[k]\,k[k] \;+\; q[k+\half]\,k[k+\half] \label{eq:A}\\
  \Bk &= q[k+\half]\,k[k] \;-\; q[k]\,k[k+\half] \label{eq:B}\\
  c_k(\delta) &= \Ak \cos\D + \Bk \sin\D \\
  &= \Rk\cos(\D - \psik), \\
  \Rk &= \operatorname{hypot}(\Ak, \Bk), \\
  \psik &= \operatorname{atan2}(\Bk, \Ak).
\end{align}
and the score is $s(\delta) = \sum_{k < \half} c_k(\delta)$. The two
coefficients are therefore fixed by the content vectors alone, and the whole of
the $\delta$-dependence sits in the single angle $\D$.

Note also what the closed form does \emph{not} depend on. $\Rk$ is a
hypot of two products of content coordinates; it is the same number at
$\delta = 1$ and at $\delta = 8192$. That is the content of
Proposition~\ref{prop:blind}, and it is what lets us speak of a per-pair
``gate'' at all. Two structural facts follow immediately, and both are exact
rather than asymptotic.

\begin{proposition}[Position-blindness]
\label{prop:blind}
The two coefficients $\Ak$ of Equation~\ref{eq:A} and $\Bk$ of
Equation~\ref{eq:B} are functions of the content vectors only, and do not
depend on $\delta$. Hence $\Rk$ and $\psik$ are position-independent, and the
pair $k$ is completely position-blind if and only if $\Bk = 0$.
\end{proposition}

\begin{proof}[Proof of Proposition~\ref{prop:blind}]
Read $\Ak$ and $\Bk$ off Equations~\ref{eq:A}--\ref{eq:B}: each is a sum or
difference of products $q[\cdot]k[\cdot]$ of content coordinates, and neither
$\delta$ nor $\invf_k$ appears. So both are $\delta$-independent, and so are
$\Rk = \operatorname{hypot}(\Ak, \Bk)$ and
$\psik = \operatorname{atan2}(\Bk, \Ak)$, which are functions of
$(\Ak, \Bk)$ alone.

For the only-if direction, note that $\Rk$ is the amplitude of the sinusoid
$c_k(\delta) = \Rk\cos(\D - \psik)$ with $\D = \delta\invf_k$ ranging over an
interval of length $|m - n|\cdot\invf_k$ as the relative distance varies. If
$\Bk = 0$ then $\psik = 0$, $c_k(\delta) = \Ak\cos\D$, which depends on
$\delta$ only through its sign on the intervals where $\cos$ changes sign; if
$\Bk \ne 0$ the pair mixes amplitude and phase and $c_k$ is non-constant
whenever $\invf_k \ne 0$. In the non-degenerate case $\invf_k > 0$ for every
$k < \half$, so $\Bk = 0$ is exactly the position-blind case, which is what
Proposition~\ref{prop:gate} uses to separate a per-pair ``gate'' from a
``phase''.
\end{proof}

\begin{proposition}[Conditionality]
\label{prop:cond}
Because $c_k(\delta)$ is a non-constant sinusoid in general, a feature's
contribution to the score is a function of $\delta$, unless every rotary pair
that feature touches happens to satisfy $\Bk = 0$.
\end{proposition}

\begin{proof}[Proof of Proposition~\ref{prop:cond}]
A feature's contribution is $\sum_k$ over the rotary pairs it touches, each term
$\Rk\cos(\delta\invf_k - \psik)$. If every such pair has $\Bk = 0$ then each
$\psik = 0$ and each term is $\Ak\cos(\delta\invf_k)$, whose dependence on
$\delta$ is confined to the sign of the cosine --- no amplitude information
about the relative distance is carried. If some touched pair has $\Bk \ne 0$,
the sum contains a genuine sinusoid in $\delta$ with nonzero frequency
$\invf_k > 0$, so it is not $\delta$-independent and in general changes sign as
$\delta$ ranges over the context window. Section~\ref{sec:cond} measures the
sign change rather than assuming it.
\end{proof}

Proposition~\ref{prop:cond} is the whole of R2, stated before any measurement.
Section~\ref{sec:cond} measures how large the effect is.

\subsection{Context extension, and what \yarn\ actually does}

\yarn\ \citep{peng2023yarn} extends context by leaving $\mathrm{base}$ alone and,
per Equation~\ref{eq:mscale}, modifying the inverse frequencies per channel: short-wavelength (high-frequency)
channels keep their $\invf$, long-wavelength (low-frequency) channels are
divided by the extension factor $s$, and the two are blended linearly in
between. It then multiplies all cosines and sines by a magnitude factor
\begin{equation}
  \mathrm{mscale} = 0.1 \ln s + 1 .
  \label{eq:mscale}
\end{equation}
Note what is \emph{not} here: $\mathrm{base}$ is unchanged, and the fastest
channel is explicitly left alone. This is the source of R3.

\paperfigure{fig01_frequency_ladder}{%
  Inverse frequency against rotary-pair index for the four schemes. RoPE is
  the reference ladder; position interpolation divides every entry by $32$;
  \yarn\ leaves the fastest $9$ of $32$ entries bit-for-bit identical to RoPE
  and ramps only the slow tail down by exactly $32\times$; the
  ``$\mathrm{base}=500000$'' ladder is a uniform reshaping. See
  Section~\ref{sec:context}.%
}{fig:ladder}

The frequency ladder of Figure~\ref{fig:ladder} makes the asymmetry explicit.
Under \yarn\ at $s = 32$, the fastest $9$ of the $32$ rotary pairs have inverse
frequencies \emph{bit-for-bit identical} to plain RoPE, and the slowest pair is
slowed by exactly a factor of $32.00$. The scheme is thus not a uniform
rescaling at all: it is a ramp that leaves $9$ of $32$ entries untouched.
Position interpolation, by contrast, differs from RoPE in all $32$ entries.

Position interpolation \citep{chen2023positionalinterp} instead divides
positions by $s$ uniformly. It is worth being precise that this is
\emph{not} the same scheme as raising $\mathrm{base}$. Dividing positions
divides every entry by $s$, whereas raising $\mathrm{base}$ to
$\mathrm{base}\cdot s$ multiplies entry $k$ by $s^{-2k/d}$:
\[
  \mathrm{inv\_freq}_k \;\mapsto\; \mathrm{base}^{-2k/d}\,s^{-1}
  \qquad\text{versus}\qquad
  \mathrm{base}^{-2k/d}\,s^{-2k/d}.
\]
The two agree only where $2k/d = 1$, that is at $k = d/2$, which lies outside
the ladder $k = 0,\dots,d/2-1$. They therefore never coincide on any channel.
The sharpest instance is the fastest channel: $\theta_0 = \mathrm{base}^{0} = 1$
for \emph{every} base, so raising the base from $10000$ to $500000$ leaves
$\theta_0$, and hence $\max_k |D_k|$, completely unchanged, while dividing
positions by $s$ divides it. This is exactly what our table shows: at
$\delta = 512$ position interpolation reaches $\max_k|D_k| = 16$ and the
$\mathrm{base} = 500000$ ladder still sits at $512$. We nonetheless measure all
four as separate rows, because they are frequently conflated.

% =====================================================================
\section{Method}

All measurements come from a from-scratch float64 NumPy implementation,
regenerated by a single entry point and written to
\texttt{results/measurements.json}. We use $d = 64$, $\mathrm{base} = 10000$,
extension factor $s = 32$, and pre-extension context $2048$, unless a table
says otherwise. Two grids are used in the structural and scheme-comparison
tables and are kept distinct throughout: a
$5$-point grid $\{1, 64, 512, 2048, 4096\}$ for the structural and
scheme-comparison tables, and a $54$-point log-spaced integer grid from $1$ to
$8192$ for the figures. Every quantity carries its grid with it, so numbers
drawn from different grids are never compared as though they were the same.

\paragraph{Measured quantities.}
\begin{itemize}
  \itemsep2pt
  \item \textbf{Structural residuals.} Max absolute error of the
    relative-position property over $42$ position pairs; max absolute error of
    $\bigl\|\R_m q\bigr\| - \bigl\|q\bigr\|$ and
    $\bigl\|\R_n k\bigr\| - \bigl\|k\bigr\|$; max relative error of bilinearity
    over $35$ score checks; max absolute error of the per-pair closed form
    against brute force.
  \item \textbf{Feature additivity.} We build $q = \sum_i f_i d_i^{(q)}$ and
    $k = \sum_j g_j d_j^{(k)}$ over $F = 8$ synthetic directions with
    non-negative coefficients, project each summand independently, and check
    $s(\delta) \stackrel{?}{=} \sum_{ij} A_{ij}(\delta)$ where
    $A_{ij} = q_i^{\!\top}\R_{\delta}k_j$.
  \item \textbf{Position-conditionality.} For each feature $i$ we record the
    signed curve $c_i(\delta)$ over the whole range, and report three summaries of
    it: whether it changes sign at all, its coefficient of variation in
    magnitude, and the residual left by replacing it with a position-independent
    constant. We deliberately do \emph{not} report
    $\max_\delta |c_i| / \min_\delta |c_i|$: that is a property of the sampling
    grid, and Section~\ref{sec:cond} explains what went wrong when we did.
  \item \textbf{Linearization.} We replace $\cos \D$ by $1 - \D^2/2$ and
    $\sin \D$ by $\D$, and measure the damage per pair, normalized by that
    pair's own amplitude $\Rk$, plus the amplitude-weighted total. The per-pair
    normalization is not a detail. The exact score is a sum of $\half$
    sinusoids in $\delta$ that cancel; as a function of $\delta$ it passes
    through zero repeatedly, so a total-score relative error
    $\bigl|\hat s - s\bigr| / \bigl|s\bigr|$ diverges at those zeros for
    reasons that have nothing to do with the quality of the approximation. Any
    scheme that reports a total-score relative error under a linearization is,
    at the least, reporting something dominated by cancellation. Normalizing by
    $\Rk$ bounds the error by $1$ in units of the pair's own contribution,
    which is a well-conditioned quantity: it is $0$ when the pair is exactly
    linear and $1$ when the error is as large as the term being approximated.
  \item \textbf{A ``linearizable'' channel} is one with $|\D| \le 0.1\rad$
    throughout the grid. This is a definition, a dimensionful threshold, not a
    measurement.
  \item \textbf{Entropy.} Attention entropy $H/\ln T$ for one query against
    $T = 256$ keys, with and without $\mathrm{mscale}$, over $18$ extension
    factors.
\end{itemize}

\paragraph{Reproducibility.} Every figure ships the array it plotted as a CSV
sidecar, so no claim in this paper is read off a pixel, and the appendix records
which file and key every quoted number came from. The audit is mechanical: each
literal in the tables is either an entry of \texttt{results/measurements.json},
an entry of a figure CSV, or a stated quotient of two such entries.

% =====================================================================
\section{Results}

\subsection{Where the non-additivity actually is}
\label{sec:softmax}

Because exactness at the score level is the point most often gotten wrong in
both directions, it is worth being precise about what the softmax does and does
not destroy.

Write the attention weight for key $n$ as
\begin{equation}
  a_n \;=\; \frac{\exp\!\bigl(s(\delta_n)/\sqrt{d}\bigr)}
  {\sum_{n'} \exp\!\bigl(s(\delta_{n'})/\sqrt{d}\bigr)} .
\end{equation}
Two separate operations happen here. The first is the map $n \mapsto
s(\delta_n)$, which by Equation~\ref{eq:additive} is exactly additive in the
features: $s(\delta_n) = \sum_{ij} f_i g_j A_{ij}(\delta_n)$, with the
$\delta_n$ fixed per key. The second is the normalization across $n$, and
\emph{that} is what produces the coupling between keys. Two keys whose feature
contributions are individually additive are not jointly additive after
division, because the denominator sums over all of them.

The practical consequence for attribution is a distinction we think is worth
making sharply. A feature's contribution to the \emph{logits} is exactly
additive and exactly position-conditional, and that is what this paper
measures. A feature's contribution to the attention \emph{weights} is neither,
and the error between the two is not small in general --- it grows as the
number of keys grows and as the logit spread grows. So an additive attribution
method applied to RoPE attention can be exactly right about the score and
still be wrong about the weights, and if the object of study is the weights,
the obstruction is the softmax rather than the position encoding.

We deliberately do not quantify that second error. Doing so would require a
choice of key set and of logit scale, and the resulting number would be a
property of that choice rather than of RoPE. It is the clearest place where
this paper stops, and it is a limit of the method rather than of the
position encoding.

\subsection{Exactness, measured}
\label{sec:exact}

% table* spans both columns: at \columnwidth the three-column body is
% ~78pt too wide and runs into the margin. Layout-only change; no content
% is added, removed or reworded.
\begin{table*}[t]
  \centering
  \caption{Structural properties of RoPE, measured. Residuals are
    max-absolute unless marked otherwise. All are at the float64 roundoff
    floor, i.e.\ none is an approximation with a meaningful residual.
    Source: \texttt{results/measurements.json} and
    \texttt{figures/fig09}\allowbreak\emph{\_}\allowbreak
    position-conditional CSV.}
  \label{tab:exact}
  \setlength{\tabcolsep}{4.5pt}
  \begin{tabular}{lrl}
    \toprule
    property & residual & source \\
    \midrule
    relative-position property & $<10^{-13}$ & JSON \\
    \quad position pairs tested & $42$ & JSON \\
    norm preservation, query & $<10^{-14}$ & JSON \\
    norm preservation, key & $<10^{-14}$ & JSON \\
    bilinearity in content & $<10^{-14}$ & JSON \\
    \quad score checks & $35$ & JSON \\
    per-pair closed form vs.\ brute force & $<10^{-14}$ & JSON \\
    \midrule
    feature additivity, $5$-pt grid & $<10^{-13}$ & JSON \\
    feature additivity, $54$-pt grid & $<10^{-13}$ & fig09 \\
    \bottomrule
  \end{tabular}
\end{table*}

Table~\ref{tab:exact} establishes the first half of R1: every residual is below
the float64 noise floor, the relative-position property over $42$ pairs, the
orthogonality behind Proposition~\ref{prop:gate} on both query and key, and
bilinearity across $35$ score checks. The per-pair closed form of
Section~\ref{sec:pair} reproduces the brute-force score to the same floor, so the
algebra above is the actual computation, not a post-hoc description of it.

We report these as ceilings rather than digits. Each is the residual of
subtracting two analytically equal expressions, so its value is whatever the host's
summation order produces: the same computation yields $3.4\times10^{-14}$ on one
platform and $3.6\times10^{-14}$ on another. Quoting more precision than the
quantity carries would be false. The claim that carries meaning is the bound, and
it is checked against the data at test time, not merely asserted here.

Two of these deserve a note on how they are tested, because the obvious test is
the wrong one. The relative-position property is \emph{not} tested by comparing
$s(i,j)$ against $s(i + c, j + c)$ with independent random content vectors at
each position; that comparison asks whether the score depends on absolute
position at all, which is a different and much stronger property, and one that
RoPE does not have. We test the correct statement --- that shifting a
\emph{fixed} pair of content vectors together by $7$ positions leaves their
score unchanged --- over $42$ pairs spanning $7$ distances and $6$ base
positions. Bilinearity is tested by rescaling one side by
$\{0, 1, 2, -3, 7.5, 10^{-3}, 10^{3}\}$ at $5$ distances, giving the $35$
score checks in the table.

A concrete instance of the closed form makes the structure visible. The
amplitude $\Rk$ of the fastest rotary pair, $k = 0$, takes the single value
$0.9526294617807264$ at every one of the $54$ distances and under both \rope\ and
\yarn\ --- one value, not fifty-four. That is
Proposition~\ref{prop:blind} visible directly in the data: the amplitude is a
property of the content vectors, and distance moves only the phase. For
contrast, a mid-frequency pair $k = 12$ has amplitude $2.3046990966686227$
under the same key and query, so amplitudes are not interchangeable across
pairs either; each pair has its own gate, and each gate is distance-free.

The additivity check is the one that matters for attribution. Writing
$q = \sum_i f_i d_i^{(q)}$ and $k = \sum_j g_j d_j^{(k)}$ and expanding both
before rotating, the score decomposes additively:
\begin{equation}
  s(\delta) \;=\; \sum_{i,j} f_i g_j
  \underbrace{q_i^{\!\top}\R_{\delta}k_j}_{A_{ij}(\delta)} .
  \label{eq:additive}
\end{equation}
Equation~\ref{eq:additive} holds to $3.6\times10^{-14}$ on the $5$-point grid
and $5.0\times10^{-14}$ on the $54$-point grid. Additive attribution to a
sparse-feature code is therefore \emph{exact} at the score level under RoPE, and
any failure mode one wishes to attribute to RoPE's position encoding is not
located here.

The consequence for the interpretation of attention is worth stating plainly,
because it is the opposite of a common claim. Attention weights are
$\operatorname{softmax}_n\, s(m,n)/\sqrt{d}$, and the softmax is where
non-additivity lives: bilinearity holds in the score and is destroyed by the
normalization across $n$. Section~\ref{sec:softmax} separates the two operations
and says what survives. We therefore do not claim that attention
\emph{output} is additively attributable. It is not, and RoPE is not why.

\subsection{Position-conditionality is the real obstruction}
\label{sec:cond}
Exact additivity does \emph{not} give position-free coefficients. In
Equation~\ref{eq:additive}, $A_{ij}(\delta)$ carries an explicit $\delta$, even
though the score it sums to depends on $\delta$ alone
(Equation~\ref{eq:relative}). By
Proposition~\ref{prop:cond}, a single feature's contribution
$c_i(\delta) = \sum_j A_{ij}(\delta)$ is a function of distance, and by
Proposition~\ref{prop:blind} it is position-free only if every pair it touches
has $\Bk = 0$.

\paperfigure{fig09_position_conditional_attribution}{%
  Per-feature contribution against relative distance. Eight synthetic
  features; each panel plots $c_i(\delta)$ against $\delta$ with the feature's
  own mean as a dotted line. Every feature's curve changes sign over the range.
  See Section~\ref{sec:cond}.%
}{fig:attrib}

Figure~\ref{fig:attrib} measures the effect directly, and the result is sharper
than a magnitude ratio. With $F = 8$ features on the $54$-point distance grid,
\textbf{every one of the eight contributions changes sign} somewhere in
$\delta \in [1, 8192]$. The magnitude of $|c_i|$ varies with a coefficient of
variation of $0.70 \pm 0.03$ across seeds, and the magnitudes at the two ends of
the range are the same size to within a factor of $10^{-0.08}$. What varies is
not how large the contribution is but which way it points.

This is the central result, and it is a stronger one than a large ratio would
have been. A feature whose attribution reverses sign cannot be summarised by any
position-independent number --- not approximately, but with the wrong sign. That
is what obstructs position-free attribution.

\begin{remark}[The ratio we previously led with was an artefact]
An earlier version of this section reported a worst-feature ratio
$\max_\delta|c_i|/\min_\delta|c_i|$ of $920\times$. We withdrew it, and it is
worth recording why, because the failure mode is general. The ratio is
$\max/\min$ of a continuous function over a sampled grid: as the grid gets
denser, $\min_\delta|c_i|$ keeps shrinking toward zero while the numerator
saturates, so the ratio grows without bound. Measured on the same vectors at
$5$, $17$, $48$, $104$, $235$ and $599$ distances it takes values spread by a
factor of $32$, and it is not even monotone in density. Across $12$ independent
seeds its standard deviation exceeds its own mean. It measures the sampling, not
the function.

The replacements above are chosen for stability rather than for size: the
coefficient of variation has $3.7\%$ seed variance, and the sign-crossing
fraction has none at all, because both are computed on a pre-registered distance
range rather than on whatever extremum a particular grid happens to land near.
All of this is recomputed in
\texttt{rope\_attribution.statistics} and written to
\texttt{results/statistics.json}; the test suite pins the artefact's presence so
it cannot be quietly reinstated.
\end{remark}

The practical reading: an additive attribution $\sum_{ij} f_i g_j A_{ij}$ with
$A_{ij}$ free of position is not merely an approximation --- it is a
mis-specification. It cannot represent the measured curves for any constant
$A_{ij}$, and the size of the failure is a measurement rather than a slogan:
replacing $c_i(\delta)$ by its per-feature mean over the range leaves a residual
some $15$ orders of magnitude above the additivity residual, and that factor
stays between $15.0$ and $15.2$ as the distance grid is made $120\times$ denser.
Unlike the withdrawn ratio, this quantity has an additive residual in its
denominator, so it does not grow without bound. The correct object
is per-feature and per-distance, which we define concretely in
Section~\ref{sec:wellposed}. Whether that object is useful in practice is a
question about downstream circuits, and is outside what we measure.

\subsection{What the well-posed object is}
\label{sec:wellposed}
It is worth being concrete about the replacement, because ``do per-distance
attribution'' is otherwise an instruction without a definition. For a query
feature $i$ and a key feature $j$ at relative distance $\delta$, define the
$2\times 2$ coefficients
\begin{equation}
  \begin{split}
    A_{ij}^{(k)}(\delta) \;=\;
    \cos\D\, A_{ij}^{A}[k] \;+\; \sin\D\, A_{ij}^{B}[k],\\
    \qquad k < \half ,
  \end{split}
\end{equation}
where $A_{ij}^{A}[k]$ and $A_{ij}^{B}[k]$ are formed from
$A_{ij} = q_i^{\!\top}k_j$ by the same split-half grouping as
Equations~\ref{eq:A}--\ref{eq:B}. Then
\begin{equation}
  \begin{split}
    A_{ij}(\delta) \;=\; \sum_{k < \half} A_{ij}^{(k)}(\delta)\\
    \;=\; \sum_{k < \half} R_k^{ij} \cos\bigl(\D - \psik^{ij}\bigr),
  \end{split}
  \label{eq:perpair}
\end{equation}
so the well-posed attribution object is not a $1\times1$ coefficient $A_{ij}$
but the $\delta$-indexed family
$\bigl\{A_{ij}^{(0)}(\delta), \dots, A_{ij}^{(\half-1)}(\delta)\bigr\}$,
together with the position-blind amplitudes $R_k^{ij}$. Equivalently, per
feature, $c_i(\delta) = \sum_j g_j A_{ij}(\delta)$ is an explicit finite sum of
$\half$ sinusoids in $\delta$ whose frequencies $\invf_k$ and amplitudes
$R_k^{ij}$ are fixed by the content.

Three properties of this object follow without further measurement. It is
\emph{finite}: $\half$ terms, no summation over a learned dictionary. It is
\emph{closed-form}: no fit, no optimization, no forward pass, since
$R_k^{ij}$ and $\psik^{ij}$ come from $q_i$ and $k_j$ by
Equations~\ref{eq:A}--\ref{eq:B}. And it is
\emph{exact}, because Equation~\ref{eq:perpair} is an identity rather than an
expansion. The cost is that it is a function, not a number, and any downstream
use must decide what to do with a function --- which is a question about
downstream circuits, and one we do not answer.

\subsection{The same identities on trained weights}
\label{sec:trained}

Everything above is measured on synthetic vectors. That is a real limitation for
a claim about what trained models do, and it is also avoidable for the narrower
question of whether the algebra survives contact with real weights. We repeat the
central measurements on activations taken from three trained checkpoints, across
$5$ relative distances and every attention head of every rotary layer.

Three findings, in decreasing order of how much they matter.

\textbf{The identities hold on trained weights, at the dtype floor.} The
relative-position property $\R_m^{\!\top}\R_n = \R_{n-m}$ is exact for real
projections, with a worst absolute error of $1.14\times10^{-12}$ across all three
checkpoints. The per-pair closed form of
Equation~\ref{eq:perpair} reproduces brute-force scores to a median relative
error of $5.06\times10^{-8}$ on \texttt{pythia-160m}, $3.57\times10^{-8}$ on
\texttt{llama-160m} and $5.85\times10^{-8}$ on \texttt{SmolLM-135M}, against a
float32 epsilon of $1.19\times10^{-7}$; bilinearity in content behaves the same
way at $8.65\times10^{-8}$, $5.99\times10^{-8}$ and $7.52\times10^{-8}$. These
are not small numbers that happen to pass: they are the numbers you get when the
implementation is correct and the arithmetic runs out of precision. A trained
checkpoint does not introduce a correction term to any of it.

\textbf{The obstruction of Section~\ref{sec:cond} is a fact about the weights.} The
share of an attention score taken by a feature's contribution is not a stable
quantity across heads. Pooling over heads, layers and prompts, the leading
feature's share ranges from $0.0108$ to $0.9999$ on \texttt{pythia-160m}, a
factor of $92.9$; on \texttt{llama-160m} it spans $0.0148$ to $0.9999$, a factor of $67.6$; on
\texttt{SmolLM-135M}, $0.0150$ to $0.9999$, a factor of $66.8$. The pooled
interquartile range alone runs from $0.117$ to $0.992$. Some heads hand
almost the entire score to one feature; others spread it. Any per-feature
attribution method that reports one number per feature has to pick a convention
for that spread, and the convention will dominate the answer.

\textbf{Most channels are position-carrying.} A channel counts as position-blind
when its share moves by less than $1\%$ across the distance range. Measured on
trained weights, the position-blind fraction is $0.0068$ on
\texttt{pythia-160m}, $0.0061$ on \texttt{llama-160m} and $0.0060$ on
\texttt{SmolLM-135M}: between $0.5\%$ and $0.7\%$, and remarkably consistent
across three unrelated architectures. The gate/phase split of
Proposition~\ref{prop:gate} is therefore not a useful approximation on real
weights --- almost everything is phase.

One negative result belongs here rather than in a footnote. We began with
\texttt{gpt2}, and it has no rotary embedding at all: position enters through a
learned absolute table of size $1024$, there is no rotary module in the
architecture, and no $\hat q/\hat k$ pair structure to attribute anything to. It
is a useful control precisely because it is the model one reaches for first.

What this section does not establish is equally worth stating. Three checkpoints
of $135$--$160$M parameters, one prompt family, and no downstream task, cannot
tell you whether per-distance attribution identifies anything a circuit cares
about. They can tell you that the algebra does not need correcting, and that the
head-to-head dispersion is enormous.

\subsection{What context extension fixes, and what it does not}
\label{sec:context}
\paperfigure{fig03_max_vs_median_angle}{%
  Maximum (solid) and median (dashed) per-channel angle against relative
  distance for all four schemes. The maximum angle is identical for RoPE and
  \yarn\ at every distance; only the median moves. See
  Section~\ref{sec:context}.%
}{fig:maxmedian}

Table~\ref{tab:spectrum} compares the four schemes on the $5$-point grid. Three
things are visible.

\paragraph{The maximum angle is not changed by \yarn.}
$\max_k|\D|$ for RoPE and \yarn\ is equal at \emph{every} distance in the grid.
At $\delta = 4096$ both are exactly $4096\rad$, and across all $54$ distances of
the finer grid the largest disagreement is $0.000\times10^{0}\rad$
(Figure~\ref{fig:maxmedian}). This is not
an accident of rounding: \yarn's ramp is constructed to leave
short-wavelength channels at their original inverse frequency, so the fastest
channel contributes exactly $1\rad$ per token under both schemes. Any
reasoning that says ``\yarn\ damps the large angles'' is wrong about the largest
one.

% table* spans both columns: six columns of numerics do not fit a
% single \columnwidth. Layout-only change; the rows and the caption are
% untouched.
\begin{table*}[t]
  \centering
  \caption{The four position-encoding schemes on the $5$-point distance grid,
    $d = 64$, $\mathrm{base} = 10^{4}$, extension factor $s = 32$, pre-extension
    context $2048$. ``Linearizable'' is the fraction of rotary pairs with
    $|\D| \le 0.1\rad$, with the exact count out of $\half = 32$ in brackets.
    ``Amp-wtd err'' is the amplitude-weighted linearization error, normalized so
    that $1$ would mean an error as large as the head's own total amplitude.
    Source: \texttt{measurements.json}, key \texttt{method\_spectrum}.}
  \label{tab:spectrum}
  \setlength{\tabcolsep}{3.6pt}
  \begin{tabular}{llrrrr}
    \toprule
    scheme & $\delta$ & $\max_k|\D|$ & $\mathrm{med}_k|\D|$
      & linearizable & amp-wtd err \\
    \midrule
    RoPE, $\mathrm{base}=10^{4}$ & 512 & 512 & 5.97
      & $0.0625\ (2/32)$ & $4.88\!\times\!10^{3}$ \\
    RoPE, $\mathrm{base}=10^{4}$ & 2048 & 2048 & 23.90
      & $0.0000\ (0/32)$ & $7.79\!\times\!10^{4}$ \\
    RoPE, $\mathrm{base}=10^{4}$ & 4096 & 4096 & 47.79
      & $0.0000\ (0/32)$ & $3.12\!\times\!10^{5}$ \\
    \midrule
    \yarn, $s=32$ & 512 & 512 & 2.67
      & $0.3438\ (11/32)$ & $4.87\!\times\!10^{3}$ \\
    \yarn, $s=32$ & 2048 & 2048 & 10.67
      & $0.2812\ (9/32)$ & $7.77\!\times\!10^{4}$ \\
    \yarn, $s=32$ & 4096 & 4096 & 21.34
      & $0.2188\ (7/32)$ & $3.11\!\times\!10^{5}$ \\
    \midrule
    interpolation & 512 & 16 & 0.19
      & $0.4375\ (14/32)$ & $5.26$ \\
    interpolation & 4096 & 128 & 1.49
      & $0.2188\ (7/32)$ & $3.08\!\times\!10^{2}$ \\
    \midrule
    $\mathrm{base}=5\!\times\!10^{5}$ & 512 & 512 & 0.91
      & $0.3438\ (11/32)$ & $2.78\!\times\!10^{3}$ \\
    $\mathrm{base}=5\!\times\!10^{5}$ & 4096 & 4096 & 7.26
      & $0.1875\ (6/32)$ & $1.77\!\times\!10^{5}$ \\
    \bottomrule
  \end{tabular}
\end{table*}

\paragraph{The median angle falls, and the linearizable fraction rises.}
$\operatorname{median}_k|\D|$ for \yarn\ is a constant $0.4464\times$ that of
RoPE at every distance tested --- $2.67$ against $5.97\rad$ at $\delta = 512$,
and $21.34$ against $47.79$ at $\delta = 4096$. A fixed ratio, not an
improvement that grows with distance, because \yarn's ramp is a fixed
per-channel rescaling and $\D$ is linear in $\delta$ for every channel.
Figure~\ref{fig:spectrum} shows the per-channel form of this: at
$\delta = 4493$ the median angle is $52.42\rad$ under RoPE against
$23.40\rad$ under \yarn, and the two distributions have visibly different
\emph{left} tails even though their right edges coincide exactly.
Consequently the fraction of channels admitting $\cos \D \mapsto 1 - \D^2/2$,
$\sin \D \mapsto \D$ rises from $0.0625$ to $0.3438$ at $\delta = 512$, and from
$0.0000$ to $0.2188$ at $\delta = 4096$. The base-raising heuristic is a
different row with different numbers again ($0.3438$ and $0.1875$), and
position interpolation is the most favourable on this metric because it damps
the fast channel too, at the cost of shrinking $\max_k|\D|$ from $4096$ to
$128\rad$.

\paperfigure{fig02_angle_spectrum}{%
  Per-channel angle $|D_k|$ against pair index at eight distances, for RoPE
  and for \yarn. The right-hand edge of both panels is identical; \yarn\ lowers
  the bulk of the spectrum. See Section~\ref{sec:context}.%
}{fig:spectrum}

\paperfigure{fig04_linearizable_fraction}{%
  Fraction of rotary pairs admitting the first-order linearization, against
  relative distance, for all four schemes. Plain RoPE reaches zero; \yarn\ and
  the other schemes retain a shrinking minority. See
  Section~\ref{sec:context}.%
}{fig:linearizable}

\paragraph{Neither rescues a global linearization.}
Figure~\ref{fig:linearizable} tracks this across the $54$-point grid. Plain
RoPE retains no linearizable channel at all from $\delta = 861$ onward, where it
is exactly $0.0000$ at every subsequent grid point, while \yarn\ retains
$11/32$ there and $4/32$ at $\delta = 8192$. So \yarn\ does not rescue the
linearization either; it lengthens the range over which a majority of channels
can be linearized. The per-channel mechanism is in
Figure~\ref{fig:pairs}: the fastest rotary pair's error reaches the size of its
own contribution at $\delta = 3$, identically under RoPE and \yarn\ --- the same
$\delta$, because \yarn\ does not touch that channel. Note how early that is.
The fastest channel turns a full revolution per token, so it is out of the
linear regime within three tokens of separation and stays out. No amount of
ramping the slow tail repairs a channel that is wrong that soon.

\paperfigure{fig05_linearization_error}{%
  Amplitude-weighted and per-pair linearization error against relative distance,
  all four schemes, log scale. The curves for RoPE and \yarn\ sit almost on top
  of one another even though their median angles differ by more than a factor of
  two: the error is dominated by the handful of fast channels both schemes leave
  untouched. See Section~\ref{sec:context}.%
}{fig:error}

\paperfigure{fig06_pair_exact_vs_linear}{%
  Exact per-pair contribution against its first-order Taylor form, for five
  selected rotary pairs, under RoPE and \yarn. The fastest pair's normalized
  error reaches $1$, meaning an error as large as the term itself, at
  $\delta = 3$ under both schemes. See Section~\ref{sec:context}.%
}{fig:pairs}

The same figure localizes the damage. At $\delta = 1$ the fastest pair $k = 0$
has an exact contribution of $-0.9525$ against a linearized value of
$-1.0607$, a normalized error of $0.1136$; the mid-frequency pair $k = 12$ has
normalized error $4.3\times10^{-6}$ at the same distance, roughly $26000$
times smaller. By $\delta = 3$ the fastest pair's normalized error is
$1.1099$ --- past $1$, meaning the approximation is now worse than simply
deleting the term --- while the slower pairs are still comfortably inside the
linear regime. This is the whole of Section~\ref{sec:context} in two numbers:
the linearization fails for the fast channels first, and \yarn does not change
which channels those are.

Figure~\ref{fig:error} is the decisive guard on this conclusion, and it is where
the optimistic reading dies. At $\delta = 8192$ the amplitude-weighted
error is $1.24572\times10^{6}$ for RoPE and $1.24179\times10^{6}$ for \yarn\,
a ratio of $0.996840$: \yarn\ buys less than $0.4\%$ on this measure. The
per-pair RMS ratio is $0.999977$, i.e.\ no meaningful change. The reason is
visible in the spectrum and follows from Proposition~\ref{prop:blind}: the error
is dominated by the few fast channels, whose amplitudes are comparable to the
slow channels' but whose angles are enormous, and \yarn\ scales those channels
by exactly $1$. Position interpolation does better on this axis ($1225.06$ at
$\delta = 8192$, roughly $1000\times$ below RoPE) precisely because it also
shrinks the fast channel --- which is the same reason it compresses the whole
spectrum into a regime a trained model does not natively use. That trade-off is
a modelling question about trained models, and we do not measure it.

\begin{remark}
The correct summary of R3 is: context extension changes \emph{which channels}
admit a linearization and \emph{when}, but does not change the fact that the
worst channel is out of range. RoPE runs out of linearizable channels first;
\yarn\ extends the range for the bulk. Neither makes the linearization of the
whole head safe, and neither is a rescue for position-free attribution --- that
failure is in R2, which no amount of angle damping addresses.
\end{remark}

\subsection{Partial RoPE gives an exactly position-free sub-score}
\label{sec:partial}
\pip\ rotates only a fraction $p = n_{\mathrm{rot}}/d$ of channels and leaves
the rest untouched. The score then splits into an unrotated part and a rotated
part, and the unrotated part is position-free by construction --- not
approximately, but exactly, because no rotation is applied to it.

\paperfigure{fig07_partial_rope}{%
  Partial RoPE at $p = 0.25$: the rotated and unrotated sub-scores against
  relative distance (left), and each sub-score's deviation from its own mean
  (right). The clean sub-score is flat to the last bit; the rotated one
  oscillates. See Section~\ref{sec:partial}.%
}{fig:partial}

Figure~\ref{fig:partial} shows the split. At $d = 64$ with
$n_{\mathrm{rot}} = 16$: the unrotated sub-score's spread over the $5$-point
grid is \emph{exactly} $0.0$, and over the finer $54$-point grid it
is also exactly $0.0$ --- the plotted constant is $0.6857345646988313$ at every
distance. The rotated sub-score, over those same grids, varies by $11.34$ and
$24.20$ respectively. The clean part carries a mean magnitude share of the
score of $0.302$ (JSON grid) and $0.187$ (figure grid), so this is a
substantially weighted, genuinely position-free bilinear term.

This is the one setting in which a position-free additive attribution is exactly
correct, and it is correct because a quarter of the channels carry no position
at all. That makes \pip\ the natural recommendation for feature-level
attribution in a model that has it.

\begin{remark}[$\pip$ preserves the norm too]
There is a subtlety here that a natural implementation gets wrong. Partial
rotary is orthogonal, and so preserves the norm exactly, as full RoPE does: we
measure a deviation below the float64 noise floor at
$n_{\mathrm{rot}} = 16$, the same as full RoPE. This holds only if the rotated
block is paired \emph{within itself} --- channel $i$ with channel
$i + n_{\mathrm{rot}}/2$, the GPT-NeoX layout used by models such as
\texttt{EleutherAI/pythia-160m}. Pairing $i$ with $i + d/2$ across the whole head
instead leaves every touched pair half rotated and half not; that variant is not
orthogonal and is not what any published model does. We mention it because the
failure is silent: the code runs, the numbers look reasonable, and only a
norm check reveals the difference. The one claim we withdraw here is any
suggestion that partial rotation damages the magnitude gate.
\end{remark}

\subsection{$\mathrm{mscale}$ is a temperature}
\label{sec:mscale}

\yarn's magnitude factor of Equation~\ref{eq:mscale} multiplies every logit,
which is precisely a temperature on the attention distribution. At $s = 32$ we
measure $\mathrm{mscale} = 1.3465736$, matching the closed form
$0.1\ln 32 + 1$ to all printed digits.

\paperfigure{fig08_mscale_entropy}{%
  Normalized attention entropy $H/\ln T$ against extension factor, with and
  without $\mathrm{mscale}$ (left), and the measured $\mathrm{mscale}$ against
  its closed form (right). $\mathrm{mscale}$ sharpens the attention
  distribution. See Section~\ref{sec:mscale}.%
}{fig:entropy}

Figure~\ref{fig:entropy} plots this. With one query against $T = 256$ keys,
$\mathrm{mscale}$ moves the normalized entropy from $H/\ln T = 0.9233$ to
$0.8651$, a drop of $0.0582$ in absolute
nats, from $5.1198$ to $4.7970$. Across the full range of $18$ extension
factors from $1$ to $128$, the effect is monotone: $H/\ln T$ falls to $0.8383$
as $\mathrm{mscale}$ rises to $1.4852$. Lowering entropy is a sharpening of the
attention distribution: fewer keys receive appreciable mass. This is a real,
measurable, and easily overlooked side effect of context extension that is not
part of its frequency story, and it is independent of everything in
Sections~\ref{sec:exact}--\ref{sec:context}. We report it as a measured
consequence and make no claim about whether it helps or hurts a trained model,
which we do not measure.

% =====================================================================
\section{Limitations}

R1, R2 and R3 are usually discussed as three separate observations, and they are
not: R1 and R2 pull in opposite directions, and R3 is orthogonal to both.
Appendix~\ref{sec:synthesis} sets out how they relate, and why context extension
is \emph{indifferent} on the axis that matters most.

In decreasing order of how much they constrain the conclusions.

\textbf{No downstream model, and no learned features.} The algebraic claims rest
on synthetic vectors with fixed seeds, and Section~\ref{sec:trained} confirms the
central identities on activations from three trained checkpoints at
$135$--$160$M parameters. Neither connects to a task: there is no retrieval
benchmark, no language-modelling evaluation and no causal patching anywhere in
this work. The $8$ directions our attribution runs over are random with
non-negative coefficients: they have the shape of a sparse code and none of its
content, so nothing here is evidence about real feature geometry or circuit
structure. Where we say \yarn\ is better on linearizable fraction, we mean better
on that number, and Section~\ref{sec:context} shows it does not track the
amplitude-weighted error. In particular, the sign reversal in
Section~\ref{sec:cond} is a property of random directions in a
$64$-dimensional head, and although Section~\ref{sec:trained} measures the
head-to-head dispersion of feature shares on real weights, we have not measured
the sign reversal itself for real features.

\textbf{Single configuration.} All measurements use $d = 64$, $\mathrm{base} =
10000$, $s = 32$, pre-extension context $2048$, and one query/key draw per
experiment. We have not swept head dimension, and the claim about which channels
dominate the error should be expected to depend on $d$ through the number of
pairs $\half$; we did not measure that dependence and do not claim it. The
entropy result of Section~\ref{sec:mscale} uses a random query against $256$
random keys, so it is a property of that synthetic configuration rather than of
any trained distribution; we claim only that $\mathrm{mscale}$ acts as a
temperature there, which is monotone across $18$ scales.

\textbf{Approximation choices.} Linearization is one particular first-order
approximation: we test $\cos \D \mapsto 1 - \D^2/2$, $\sin \D \mapsto \D$ with
a $0.1\rad$ threshold. A different truncation, or a per-channel threshold, would
move the linearizable fractions, and would leave the amplitude-weighted
conclusion intact only if it did not change which channels dominate. We verified
that qualitatively through the per-pair error but did not sweep thresholds.

\textbf{Scope of the exactness claim.} Exactness at the score level is a
statement about a bilinear map, and its practical reach is bounded by three
things we have not crossed. It does not survive the softmax: attention
\emph{weights} are not additive, so a feature's share of attention mass is not
its share of the score, and we do not measure the discrepancy. It does not
survive multi-head composition, where the position-conditional structure we
describe is per-head and no head shares our synthetic projections. And it does
not license a claim about attribution methods we did not run; we take
attribution patching as a reference point, not as something we reproduce. A
reader should treat exactness as removing one stated objection to RoPE-based
attribution, not as validating the method. Relatedly, Section~\ref{sec:cond}
shows that a position-free scalar $A_{ij}$ is mis-specified; it does not show
that per-distance attribution is practically useful, which would require
circuits, patching, and a trained model.

% =====================================================================
\section{Related Work}

\textbf{Position encodings and length extrapolation.} RoPE
\citep{su2021ropeformer} is the scheme studied here. \alpibib\
\citep{press2021train} is an instructive contrast: it removes position from the
attention pattern via a linear bias, so it has no per-channel angle spectrum and
no version of the linearization problem of Section~\ref{sec:context}. On
extension, position interpolation \citep{chen2023positionalinterp} rescales
positions uniformly --- the scheme measured as \texttt{position\_interpolation},
which our table shows is geometrically the most favourable on linearizable
fraction. \yarn\ \citep{peng2023yarn} keeps the base and ramps per frequency,
which is what makes $\max_k|\D|$ invariant under it; to our knowledge this is a
point on which the framing of ``\yarn\ damps large angles'' is easily misread.
\lengthxtrap\ \citep{sun2022lengthextrap} covers the partial case of
Section~\ref{sec:partial}, \longrope\ \citep{ding2024longrope} per-channel
non-uniform extension, and \ropetheory\ \citep{barbero2024roundandround} RoPE's
frequency content from an optimization perspective; we neither compare with nor
reproduce the last. Finally, the ``raise $\mathrm{base}$ from $10000$ to
$500000$'' recipe is not \yarn, and is also not position interpolation:
raising the base multiplies channel $k$ by $s^{-2k/d}$ rather than by $s^{-1}$,
so it leaves $\theta_0 = 1$ and hence $\max_k|\D|$ untouched. We measure it
separately as \texttt{base\_500k\_legacy\_claim} for that reason.

\textbf{Feature-level attribution.} The additive decomposition of activations
into learned features originates in the superposition line
\citep{elhage2022toymodels}, with sparse autoencoders as an operational extractor
\citep{cunningham2023sparseautoencoders}. Attribution on attention heads
includes attribution patching \citep{syed2023attributionpatching}, which we take
as the methodological reference for ``which features contributed to this
computation'', and the tuned lens \citep{belrose2023tunedlens}. Feature-level
causal graphs appear in sparse feature circuits \citep{marks2024sparsefeaturecircuits}
and transcoder-based discovery \citep{dunefsky2024transcoders}. What is common
to all of them is that they operate at a \emph{fixed} position, or integrate
over positions treating each as an independent example: the positional encoding
is part of the forward pass, not the attribution target.

\textbf{Our relationship to that literature.} Those methods ask which features
caused a behaviour in a trained model. We ask a strictly prior question about
the position encoding alone: given a RoPE head, what is the exact algebraic form
of a feature's contribution, and is it a scalar? Our answer --- no, it is a
function of distance --- is a structural precondition that any of those methods
must handle and that, to our knowledge, has not been stated in this form. It is
not a competing attribution method and it evaluates none. We mention the relation
because it is a limitation of this paper rather than of the field: our
per-distance attribution has not been connected to a circuit, and we have not
shown that measuring it identifies anything about a trained model.

% =====================================================================
\section{Conclusion}

RoPE's effect on feature-level attention attribution is exactly characterizable
in the relative frame, where the score is $q^{\!\top}\R_{\delta}k$ and each of
the $\half$ rotary pairs contributes a sinusoid with a position-independent
amplitude and a linear phase. Three consequences. Attribution at the score level
is exact, to $3.6\times10^{-14}$, and orthogonality makes the magnitude gate
exactly position-independent, so the gate/phase split is a theorem rather than
an approximation. Position-conditionality, not bilinearity, is the operative
obstruction: a feature's contribution reverses sign across
distance while the total stays additive to $5.0\times10^{-14}$, so
position-free per-feature attribution is mis-specified under RoPE and the
per-feature, per-distance one is the correct object. And context extension
changes the range over which channels can be linearized without changing the
worst channel at all --- $\max_k|\D|$ is identical for RoPE and \yarn\ at every
distance tested --- so \yarn\ buys a larger linearizable majority, at an
amplitude-weighted error ratio of only $0.996840$, not a safe global
linearization. Partial RoPE is the one construction that delivers an exactly
position-free sub-score, at the cost of losing norm preservation.
Section~\ref{sec:synthesis} collects how the three results relate, and the
short version is that exactness, position-conditionality, and linearizability
are three separate problems that are easy to conflate.

The measurements are structural and model-independent by design, and we have
been explicit about what that excludes: no trained model, no learned features,
no retrieval benchmark, no accuracy figure. What is claimed here is a correct
and carefully measured characterization of the score's structure under RoPE, and
nothing beyond it.

\subsection*{A note on two corrections}

Two statements are common enough in informal write-ups of this area that we
record what we measured against them, since both are cases where the natural
explanation is wrong in an instructive way. The first is that ``RoPE breaks
bilinearity''. It does not, at the score level: doubling the query doubles the
score to a relative error of $7.4\times10^{-16}$. Any observed super-linearity
comes from changing what is being measured --- typically comparing a score
against an attention output, which involves the softmax. The second is that
``YaRN works by raising the base from $10000$ to $500000$''. It does not.
\yarn\ keeps the base and ramps per frequency, leaving the fastest $9$ of $32$
pairs bit-for-bit unchanged; the base-raising recipe is a distinct scheme again,
since it rescales channel $k$ by $s^{-2k/d}$ rather than by $s^{-1}$, and
appears as its own row in
Table~\ref{tab:spectrum}. Both corrections push the same way: toward
disentangling the position encoding from the operations people attribute to it.

% =====================================================================
\appendix
\section{Provenance of every number}

Every quantitative claim in this paper is either a literal in the table below,
or a value derived from it by an arithmetic operation stated where it is used.
No number was typed from memory.

\paragraph{Configuration.} $d = 64$, $\mathrm{base} = 10000$, extension factor
$s = 32$, pre-extension context $2048$, $T = 256$ keys; from
\texttt{measurements.json} keys \texttt{structural\_facts.dim},
\texttt{structural\_facts.base}, \texttt{mscale\_entropy.scale},
\texttt{mscale\_entropy.n\_keys}, and the shared-configuration section of
\texttt{figures/README.md}.

\paragraph{Claims and their source files.}
% Full-width float: this three-column audit table is ~216pt wider than a
% single \columnwidth. Layout-only; every row is unchanged.
\begin{table*}[t]
\centering
\caption{Provenance of every literal quoted in the paper.}
\label{tab:provenance}
\footnotesize
\setlength{\tabcolsep}{1.9pt}
\begin{tabular}{lll}
\toprule
paper location & value(s) & file and key \\
\midrule
Tab.~\ref{tab:exact} & all rows & \texttt{measurements.json :: structural\_facts} \\
 & additivity rows & \texttt{feature\_attribution}; \texttt{fig09} CSV \\
Sec.~\ref{sec:cond} & $8/8$ features cross zero; $\mathrm{CV}=0.70$, drift $0$ & \texttt{results/statistics.json} \\
  & $15$ orders, bracket $15.0$-$15.2$ & \texttt{results/statistics.json :: position\_free\_error} \\
Sec.~\ref{sec:context} & Tab.~\ref{tab:spectrum} & \texttt{measurements.json :: method\_spectrum} \\
 & $0.000\rad$ agreement & \texttt{fig03} CSV, \texttt{max\_abs\_angle\_rad} \\
 & $0.4464$, $1.0000$ & ratios of fig03 medians and maxima \\
 & $\delta = 861$, $4/32$ & \texttt{fig04} CSV \\
 & $1.24572\!\times\!10^{6}$ etc. & \texttt{fig05} CSV \\
 & $1225.06$, $\delta = 3$ & \texttt{fig05}, \texttt{fig06} CSVs \\
Sec.~\ref{sec:context} & $9$ of $32$, $32.00\times$ & \texttt{fig01} CSV, unchanged-pair count \\
 & $52.42$, $23.40$ & \texttt{fig02} CSV, \texttt{abs\_angle\_rad} \\
Sec.~\ref{sec:exact} & $0.9526294618$, $2.3046990967$ & \texttt{fig06} CSV, \texttt{amplitude\_R\_k} \\
 & $-0.9525$, $-1.0607$, $0.1136$, $4.3\!\times\!10^{-6}$ & \texttt{fig06} CSV \\
Sec.~\ref{sec:background} & $9$ of $32$ bit-for-bit & \texttt{fig01} CSV, column 3 vs column 1 \\
Sec.~\ref{sec:partial} & $0.0$, $11.34$, $24.20$ & \texttt{measurements.json :: partial\_rope}; \texttt{fig07} CSV \\
 & $0.302$, $0.187$ & \texttt{clean\_magnitude\_share\_mean}; fig07 \texttt{clean\_subscore} \\
 & $<10^{-13}$, $1.776\!\times\!10^{-15}$ & \texttt{partial\_norm\_deviation}; full-RoPE floor \\
Sec.~\ref{sec:mscale} & $1.3465736$ & \texttt{measurements.json :: mscale\_entropy.mscale} \\
 & $0.9233 \to 0.8651$ & \texttt{entropy\_ratio\_unscaled}, \texttt{\_mscaled} \\
 & $0.8383$, $1.4852$ & \texttt{fig08} CSV at $\mathrm{scale} = 128$ \\
\bottomrule
\end{tabular}
\end{table*}

\paragraph{Derived quantities and their derivations.} The fixed median-angle
ratio quoted above is a ratio of medians at fixed $\delta$; it is constant across
distances because \yarn's ramp is a fixed per-channel rescaling and $\D$ is
linear in $\delta$ for every channel. Fractions converted to counts out of $32$
are exact, since the denominator is $\half = 32$ pairs and every fraction is a
multiple of $1/32$. Error ratios are quotients of two entries of the same CSV at
the same $\delta$. The sign-crossing fraction, the coefficient of variation and
the grid-independence drift quoted in Section~\ref{sec:cond} come from
\texttt{results/statistics.json} and are recomputed by
\texttt{rope\_attribution.statistics}; the withdrawn ratio is recomputed there
too, from six grid densities, so that its instability is part of the record
rather than a claim about a former draft. The ``less
than $0.4\%$'' claim in Section~\ref{sec:context} is $1 - 0.996840$. The
``roughly $1000\times$'' comparison divides
$1.24572\times10^{6}$ by $1225.06$.

\paragraph{Exact identities reported as equalities.} The claims that
$\max_k|\D|$ is identical for RoPE and \yarn, that $\Rk$ is a single value
across all distances, that the clean partial-RoPE sub-score has spread
$0.0$, and that the table fractions are multiples of $1/32$, are checked as
exact float equalities, not to a tolerance; the residue in each case is
$0.0$ or a single ULP, as recorded above.

\paragraph{Citations.} Every cited entry was verified against the arXiv API by
identifier or by title search; the identifier, exact title and author list in
\texttt{references.bib} are as returned by that API. No citation in this file
was included from memory.

\bibliographystyle{plainnat}
\bibliography{references}


% Moved out of the main text to respect the nine-page limit.
% Content unchanged; only its placement.

\subsection{The three results are not independent}
\label{sec:synthesis}

It is worth stating explicitly how R1, R2 and R3 relate, because they are
usually discussed as three separate observations and they are not.

R1 and R2 pull in opposite directions and both are correct. The additive
decomposition of Equation~\ref{eq:additive} is exact to
$5.0\times10^{-14}$ --- the coefficients $A_{ij}$ exist and are real numbers,
not fudges. But those real numbers are functions of $\delta$, and they vary by
over the distance range. Exactness therefore buys nothing
against position-conditionality: the map from features to score is exactly
additive at every fixed distance and wildly non-constant across distances. This
is the sense in which position-conditionality, not bilinearity, is the operative
obstruction, and it is what separates R1 from R2: an attribution scheme that
computes $A_{ij}$ correctly at $\delta = 8$ and reuses it at $\delta = 512$ has
made an error far larger than any it would make by ignoring RoPE altogether.

R3 is orthogonal to both, and this is the part most likely to be
misunderstood. The linearization failure is a statement about the
\emph{approximation} $\cos \D \mapsto 1 - \D^2/2$, not about attribution. One
could perfectly well attribute a RoPE head exactly --- that is R1 --- and still
be unable to linearize it. Conversely, one could linearize a head perfectly and
still have no position-free attribution, which is the situation under full
\rope\ at any $\delta > 3$. So the two literatures --- ``we need a wider
context'' and ``we need position-free feature attribution'' --- are addressing
different problems, and a method that answers one does not answer the other.
What \yarn\ does is extend the distance range over which a majority of channels
can be linearized, from $0/32$ to $11/32$ at $\delta \approx 861$. What it does
not do is make the worst channel linearizable, and the worst channel is what
determines whether the head as a whole can be treated as linear.

We mention this because the natural reading of Table~\ref{tab:spectrum} is
that \yarn\ is ``better'', and it is better on one axis and indifferent on the
other. The honest summary is the one in the remark above: context extension
changes which channels admit a linearization and when, and does not change the
fact that the worst channel is out of range.

\end{document}
